\subsection{Model Derivation and Signal Definitions}
\begin{table}[H]\centering\small
\caption{Model notation and units.}\label{tab:signals}
\begin{tabularx}{\linewidth}{lXl}\toprule\thead
Symbol & Description & Unit\\\midrule
$x,y,v=\dot y$ & Reference altitude, altitude, vertical velocity & m, m, m/s\\
$f_t,g,d_a$ & Mass-normalised thrust, gravity, disturbance & m/s$^2$\\
$b,k_T$ & Drag coefficient; local voltage-to-thrust slope & s$^{-1}$; m/(s$^2$ V)\\
$u,u_h,\delta u$ & Command, hover command, $u-u_h$ & V\\
$A,\tau$ & Velocity gain; velocity time constant & m/(s V); s\\
$k_s$ & Position-to-voltage coefficient & V/m\\
$r_v,y_v,e_v$ & Reference, sensor, error voltage & V\\
$K_p$ & Proportional voltage-error gain & V/V\\\bottomrule
\end{tabularx}\end{table}
Start from $\ddot y=f_t-b\dot y-g+d_a$, with $f_t=g+k_T(u-u_h)$. For $d_a=0$, zero initial conditions and $v=\dot y$,
\begin{align}
\dot v+bv&=k_T\delta u,\\
(s+b)V(s)&=k_T\Delta U(s),\\
G_v(s)&=\frac{k_T}{s+b}=\frac{k_T/b}{s/b+1}=\frac{A}{\tau s+1},\\
G_y(s)&=\frac{G_v(s)}s=\frac{A}{s(\tau s+1)}.
\end{align}
The identified parameters give $b=1/8.1=0.123456790$ s$^{-1}$ and $k_T=14.2748/8.1=1.762320988$ m/(s$^2$ V). The hover voltage does not imply $k_T=g/u_h$.

For negative feedback, $E_y=X-Y$ and $L=K_pk_sG_y$:
\begin{equation}
\frac{E_y(s)}{X(s)}=\frac1{1+L(s)}=\frac{\tau s^2+s}{\tau s^2+s+AK_pk_s}.
\end{equation}
For a step $X=x_0/s$, the final-value theorem gives $e_\infty=\lim_{s\to0}sE_y(s)=0$. For a ramp $X=a_r/s^2$,
\begin{equation}
e_\infty=\frac{a_r}{AK_pk_s},\qquad K_v=AK_pk_s=0.190106076.
\end{equation}
A parabolic reference has unbounded steady error. A constant acceleration disturbance $d_0$ instead gives $y_\infty-x_\infty=d_0/(k_TK_pk_s)$. For an input bias, $d_0=k_T\delta u_b$, so $\delta y=\delta u_b/(K_pk_s)$.
